\documentclass[10pt,a4paper]{amsart}
\usepackage{amsmath,amssymb,amsthm,graphicx}
% Retained minimal preamble for the Stacks Project chapter varieties.tex.
% Verbatim excerpt of the chapter preamble: Xy-pic setup (native preamble lines 14-18), the native theorem
% declarations with the mathematical macros (native preamble lines 154-217), and the project own no-op
% definitions of the reference/slogan/history display environments (native wrapper.tex lines 1-3) that the
% counted statement uses for its slogan. Copied mechanically; nothing added, removed or redefined.
% For commutative diagrams we use Xy-pic
\usepackage[all]{xy}

% We use 2cell for 2-commutative diagrams.
\xyoption{2cell}
% Theorem environments.
%
\theoremstyle{plain}
\newtheorem{theorem}[subsection]{Theorem}
\newtheorem{proposition}[subsection]{Proposition}
\newtheorem{lemma}[subsection]{Lemma}

\theoremstyle{definition}
\newtheorem{definition}[subsection]{Definition}
\newtheorem{example}[subsection]{Example}
\newtheorem{exercise}[subsection]{Exercise}
\newtheorem{situation}[subsection]{Situation}

\theoremstyle{remark}
\newtheorem{remark}[subsection]{Remark}
\newtheorem{remarks}[subsection]{Remarks}

\numberwithin{equation}{subsection}

% Macros
%
\def\lim{\mathop{\mathrm{lim}}\nolimits}
\def\colim{\mathop{\mathrm{colim}}\nolimits}
\def\Spec{\mathop{\mathrm{Spec}}}
\def\Hom{\mathop{\mathrm{Hom}}\nolimits}
\def\Ext{\mathop{\mathrm{Ext}}\nolimits}
\def\SheafHom{\mathop{\mathcal{H}\!\mathit{om}}\nolimits}
\def\SheafExt{\mathop{\mathcal{E}\!\mathit{xt}}\nolimits}
\def\Sch{\mathit{Sch}}
\def\Mor{\mathop{\mathrm{Mor}}\nolimits}
\def\Ob{\mathop{\mathrm{Ob}}\nolimits}
\def\Sh{\mathop{\mathit{Sh}}\nolimits}
\def\NL{\mathop{N\!L}\nolimits}
\def\CH{\mathop{\mathrm{CH}}\nolimits}
\def\proetale{{pro\text{-}\acute{e}tale}}
\def\etale{{\acute{e}tale}}
\def\QCoh{\mathit{QCoh}}
\def\Ker{\mathop{\mathrm{Ker}}}
\def\Im{\mathop{\mathrm{Im}}}
\def\Coker{\mathop{\mathrm{Coker}}}
\def\Coim{\mathop{\mathrm{Coim}}}

% Boxtimes
%
\DeclareMathSymbol{\boxtimes}{\mathbin}{AMSa}{"02}

%
% Macros for moduli stacks/spaces
%
\def\QCohstack{\mathcal{QC}\!\mathit{oh}}
\def\Cohstack{\mathcal{C}\!\mathit{oh}}
\def\Spacesstack{\mathcal{S}\!\mathit{paces}}
\def\Quotfunctor{\mathrm{Quot}}
\def\Hilbfunctor{\mathrm{Hilb}}
\def\Curvesstack{\mathcal{C}\!\mathit{urves}}
\def\Polarizedstack{\mathcal{P}\!\mathit{olarized}}
\def\Complexesstack{\mathcal{C}\!\mathit{omplexes}}
% \Pic is the operator that assigns to X its picard group, usage \Pic(X)
% \Picardstack_{X/B} denotes the Picard stack of X over B
% \Picardfunctor_{X/B} denotes the Picard functor of X over B
\def\Pic{\mathop{\mathrm{Pic}}\nolimits}
\def\Picardstack{\mathcal{P}\!\mathit{ic}}
\def\Picardfunctor{\mathrm{Pic}}
\def\Deformationcategory{\mathcal{D}\!\mathit{ef}}
\newenvironment{reference}{}{}
\newenvironment{slogan}{}{}
\newenvironment{history}{}{}
\usepackage{cleveref}
\crefname{originallemma}{original lemma}{original lemmas}
\Crefname{originallemma}{Original lemma}{Original lemmas}
\makeatletter
\newlabel{lemma-geometrically-reduced-p}{{0CCE}{source}{}{sourcepointer-1}{}}
\newlabel{lemma-geometrically-reduced-p@cref}{{[originallemma][0][]0CCE}{[1][0][]source}{}{}{}}
\makeatother
\crefname{originallemma}{original lemma}{original lemmas}
\Crefname{originallemma}{Original lemma}{Original lemmas}
\makeatletter
\newlabel{lemma-geometrically-reduced-dense-smooth-open}{{056V}{source}{}{sourcepointer-2}{}}
\newlabel{lemma-geometrically-reduced-dense-smooth-open@cref}{{[originallemma][0][]056V}{[1][0][]source}{}{}{}}
\makeatother
\crefname{originallemma}{original lemma}{original lemmas}
\Crefname{originallemma}{Original lemma}{Original lemmas}
\makeatletter
\newlabel{lemma-geometrically-regular-smooth}{{038X}{source}{}{sourcepointer-3}{}}
\newlabel{lemma-geometrically-regular-smooth@cref}{{[originallemma][0][]038X}{[1][0][]source}{}{}{}}
\makeatother
\crefname{originallemma}{original lemma}{original lemmas}
\Crefname{originallemma}{Original lemma}{Original lemmas}
\makeatletter
\newlabel{lemma-generic-points-geometrically-reduced}{{04KS}{source}{}{sourcepointer-4}{}}
\newlabel{lemma-generic-points-geometrically-reduced@cref}{{[originallemma][0][]04KS}{[1][0][]source}{}{}{}}
\makeatother
\crefname{originallemma}{original lemma}{original lemmas}
\Crefname{originallemma}{Original lemma}{Original lemmas}
\makeatletter
\newlabel{lemma-relative-frobenius}{{0CCB}{source}{}{sourcepointer-5}{}}
\newlabel{lemma-relative-frobenius@cref}{{[originallemma][0][]0CCB}{[1][0][]source}{}{}{}}
\makeatother
\crefname{originallemma}{original lemma}{original lemmas}
\Crefname{originallemma}{Original lemma}{Original lemmas}
\makeatletter
\newlabel{lemma-relative-frobenius-finite}{{0CCD}{source}{}{sourcepointer-6}{}}
\newlabel{lemma-relative-frobenius-finite@cref}{{[originallemma][0][]0CCD}{[1][0][]source}{}{}{}}
\makeatother
\crefname{originalreference}{original reference}{original references}
\Crefname{originalreference}{Original reference}{Original references}
\makeatletter
\newlabel{morphisms-definition-degree}{{02NY}{source}{}{sourcepointer-7}{}}
\newlabel{morphisms-definition-degree@cref}{{[originalreference][0][]02NY}{[1][0][]source}{}{}{}}
\makeatother
\crefname{originallemma}{original lemma}{original lemmas}
\Crefname{originallemma}{Original lemma}{Original lemmas}
\makeatletter
\newlabel{lemma-relative-frobenius-endomorphism-identity}{{0CCA}{source}{}{sourcepointer-8}{}}
\newlabel{lemma-relative-frobenius-endomorphism-identity@cref}{{[originallemma][0][]0CCA}{[1][0][]source}{}{}{}}
\makeatother
\crefname{originalreference}{original reference}{original references}
\Crefname{originalreference}{Original reference}{Original references}
\makeatletter
\newlabel{morphisms-lemma-smooth-etale-over-affine-space}{{054L}{source}{}{sourcepointer-9}{}}
\newlabel{morphisms-lemma-smooth-etale-over-affine-space@cref}{{[originalreference][0][]054L}{[1][0][]source}{}{}{}}
\makeatother
\crefname{originalreference}{original reference}{original references}
\Crefname{originalreference}{Original reference}{Original references}
\makeatletter
\newlabel{morphisms-lemma-composition-relative-dimension-d}{{02NL}{source}{}{sourcepointer-10}{}}
\newlabel{morphisms-lemma-composition-relative-dimension-d@cref}{{[originalreference][0][]02NL}{[1][0][]source}{}{}{}}
\makeatother
\crefname{originalreference}{original reference}{original references}
\Crefname{originalreference}{Original reference}{Original references}
\makeatletter
\newlabel{morphisms-lemma-degree-composition}{{02NZ}{source}{}{sourcepointer-11}{}}
\newlabel{morphisms-lemma-degree-composition@cref}{{[originalreference][0][]02NZ}{[1][0][]source}{}{}{}}
\makeatother
\title{Reducedness of the Frobenius twist, geometric reducedness and existence of a smooth open are equivalent, and then the relative Frobenius is finite dominant of degree p to the dimension}
\author{The Stacks Project Authors}
\date{stacks-project:varieties.tex:local-snapshot-20261011; GNU FDL 1.2 or later}
\begin{document}
\maketitle
\noindent\textit{Editorial note.} All delivered mathematical text is the author's own, copied verbatim from the pinned Stacks Project chapter file varieties.tex; no statement or proof step was rewritten, reordered or removed.\par\medskip
\noindent\textit{Editorial note.} The counted span cites twelve cross-chapter labels, each recorded in assembly.external\_references with its canonical tag as a navigation pointer.\par\medskip
\noindent\textit{Editorial note.} The proof cites no bibliography item inside the counted span.\par\medskip
\noindent\textit{Allowed standard prerequisites (editorial).}\begin{itemize}
\item Finiteness of the relative Frobenius. Stacks Project, Varieties, Lemma lemma-relative-frobenius-finite (tag 0CC9). Makes the relative Frobenius a finite morphism, so that its degree is defined.
\item Equivalent characterizations of geometric reducedness. Stacks Project, Varieties, Lemmas lemma-geometrically-reduced-p (tag 0CC6), lemma-geometrically-reduced-dense-smooth-open (tag 0CCH) and lemma-generic-points-geometrically-reduced (tag 0CCE). Supply the three implications making up the equivalence.
\end{itemize}\medskip
\ifdefined\hypertarget\hypertarget{sourcepointer-1}{}\fi
\noindent\textit{Source reference.} Original lemma 0CCE: sources/stacks-chow/tags.txt, source lines 15504--15504. Cross-chapter standard prerequisite cited by the author inside the counted span; it is used as a pointer and its own proof is not part of this item.\par\medskip
\ifdefined\hypertarget\hypertarget{sourcepointer-2}{}\fi
\noindent\textit{Source reference.} Original lemma 056V: sources/stacks-chow/tags.txt, source lines 6723--6723. Cross-chapter standard prerequisite cited by the author inside the counted span; it is used as a pointer and its own proof is not part of this item.\par\medskip
\ifdefined\hypertarget\hypertarget{sourcepointer-3}{}\fi
\noindent\textit{Source reference.} Original lemma 038X: sources/stacks-chow/tags.txt, source lines 4252--4252. Cross-chapter standard prerequisite cited by the author inside the counted span; it is used as a pointer and its own proof is not part of this item.\par\medskip
\ifdefined\hypertarget\hypertarget{sourcepointer-4}{}\fi
\noindent\textit{Source reference.} Original lemma 04KS: sources/stacks-chow/tags.txt, source lines 5983--5983. Cross-chapter standard prerequisite cited by the author inside the counted span; it is used as a pointer and its own proof is not part of this item.\par\medskip
\ifdefined\hypertarget\hypertarget{sourcepointer-5}{}\fi
\noindent\textit{Source reference.} Original lemma 0CCB: sources/stacks-chow/tags.txt, source lines 15501--15501. Cross-chapter standard prerequisite cited by the author inside the counted span; it is used as a pointer and its own proof is not part of this item.\par\medskip
\ifdefined\hypertarget\hypertarget{sourcepointer-6}{}\fi
\noindent\textit{Source reference.} Original lemma 0CCD: sources/stacks-chow/tags.txt, source lines 15503--15503. Cross-chapter standard prerequisite cited by the author inside the counted span; it is used as a pointer and its own proof is not part of this item.\par\medskip
\ifdefined\hypertarget\hypertarget{sourcepointer-7}{}\fi
\noindent\textit{Source reference.} Original reference 02NY: sources/stacks-chow/tags.txt, source lines 3497--3497. Cross-chapter standard prerequisite cited by the author inside the counted span; it is used as a pointer and its own proof is not part of this item.\par\medskip
\ifdefined\hypertarget\hypertarget{sourcepointer-8}{}\fi
\noindent\textit{Source reference.} Original lemma 0CCA: sources/stacks-chow/tags.txt, source lines 15500--15500. Cross-chapter standard prerequisite cited by the author inside the counted span; it is used as a pointer and its own proof is not part of this item.\par\medskip
\ifdefined\hypertarget\hypertarget{sourcepointer-9}{}\fi
\noindent\textit{Source reference.} Original reference 054L: sources/stacks-chow/tags.txt, source lines 6644--6644. Cross-chapter standard prerequisite cited by the author inside the counted span; it is used as a pointer and its own proof is not part of this item.\par\medskip
\ifdefined\hypertarget\hypertarget{sourcepointer-10}{}\fi
\noindent\textit{Source reference.} Original reference 02NL: sources/stacks-chow/tags.txt, source lines 3484--3484. Cross-chapter standard prerequisite cited by the author inside the counted span; it is used as a pointer and its own proof is not part of this item.\par\medskip
\ifdefined\hypertarget\hypertarget{sourcepointer-11}{}\fi
\noindent\textit{Source reference.} Original reference 02NZ: sources/stacks-chow/tags.txt, source lines 3498--3498. Cross-chapter standard prerequisite cited by the author inside the counted span; it is used as a pointer and its own proof is not part of this item.\par\medskip
\begin{lemma}
\label{lemma-inseparable-deg-p-smooth}
Let $k$ be a field of characteristic $p > 0$. Let $X$ be a variety over $k$.
The following are equivalent
\begin{enumerate}
\item $X^{(p)}$ is reduced,
\item $X$ is geometrically reduced,
\item there is a nonempty open $U \subset X$ smooth over $k$.
\end{enumerate}
In this case $X^{(p)}$ is a variety over $k$ and $F_{X/k} : X \to X^{(p)}$
is a finite dominant morphism of degree $p^{\dim(X)}$.
\end{lemma}

\begin{proof}
We have seen the equivalence of (1) and (2) in
Lemma \ref{lemma-geometrically-reduced-p}.
We have seen that (2) implies (3) in
Lemma \ref{lemma-geometrically-reduced-dense-smooth-open}.
If (3) holds, then $U$ is geometrically reduced
(see for example Lemma \ref{lemma-geometrically-regular-smooth})
and hence $X$ is geometrically reduced by
Lemma \ref{lemma-generic-points-geometrically-reduced}.
In this way we see that (1), (2), and (3) are equivalent.

\medskip\noindent
Assume (1), (2), and (3) hold. Since $F_{X/k}$ is a homeomorphism
(Lemma \ref{lemma-relative-frobenius}) we see that $X^{(p)}$ is a variety.
Then $F_{X/k}$ is finite by Lemma \ref{lemma-relative-frobenius-finite}.
It is dominant as it is surjective. To compute the degree
(Morphisms, Definition \ref{morphisms-definition-degree})
it suffices to compute the degree of $F_{U/k} : U \to U^{(p)}$
(as $F_{U/k} = F_{X/k}|_U$ by
Lemma \ref{lemma-relative-frobenius-endomorphism-identity}).
After shrinking $U$ a bit we may assume there exists an
\'etale morphism $h : U \to \mathbf{A}^n_k$, see
Morphisms, Lemma \ref{morphisms-lemma-smooth-etale-over-affine-space}.
Of course $n = \dim(U)$ because
$\mathbf{A}^n_k \to \Spec(k)$ is smooth of relative dimension $n$,
the \'etale morphism $h$ is smooth of relative dimension $0$, and
$U \to \Spec(k)$ is smooth of relative dimension $\dim(U)$
and relative dimensions add up correctly
(Morphisms, Lemma \ref{morphisms-lemma-composition-relative-dimension-d}).
Observe that $h$ is a generically finite dominant morphism
of varieties, and hence $\deg(h)$ is defined.
By Lemma \ref{lemma-relative-frobenius-endomorphism-identity}
we have a commutative diagram
$$
\xymatrix{
U \ar[rr]_{F_{X/k}} \ar[d]_h & & U^{(p)} \ar[d]^{h^{(p)}} \\
\mathbf{A}^n_k \ar[rr]^{F_{\mathbf{A}^n_k/k}} & &
(\mathbf{A}^n_k)^{(p)}
}
$$
Since $h^{(p)}$ is a base change of $h$ it is \'etale as well
and it follows that $h^{(p)}$ is a generically finite dominant
morphism of varieties as well. The degree of $h^{(p)}$ is the
degree of the extension $k(X^{(p)})/k((\mathbf{A}^n_k)^{(p)})$
which is the same as the degree of the extension $k(X)/k(\mathbf{A}^n_k)$
because $h^{(p)}$ is the base change of $h$ (small detail omitted).
By multiplicativity of degrees
(Morphisms, Lemma \ref{morphisms-lemma-degree-composition})
it suffices to show that the degree of $F_{\mathbf{A}^n_k/k}$
is $p^n$. To see this observe that
$(\mathbf{A}^n_k)^{(p)} = \mathbf{A}^n_k$ and that
$F_{\mathbf{A}^n_k/k}$ is given by the map sending the
coordinates to their $p$th powers.
\end{proof}

\end{document}
